\documentclass[11pt,a4paper]{article}

\usepackage[utf8]{inputenc}
\usepackage[T1]{fontenc}
\usepackage[norsk]{babel}
\usepackage{lmodern}
\usepackage{microtype}
\usepackage{geometry}
\usepackage{amsmath,amssymb}
\usepackage{booktabs}
\usepackage{tabularx}
\usepackage{array}
\usepackage{enumitem}
\usepackage{hyperref}
\usepackage{xcolor}

\geometry{margin=2.5cm}
\setlength{\parindent}{0pt}
\setlength{\parskip}{0.6em}
\renewcommand{\arraystretch}{1.15}

\title{Teknisk dokumentasjon -- modeller og estimering}
\author{}
\date{}

\begin{document}

\maketitle

Dette dokumentet beskriver de statistiske modellene som brukes til å
framskrive \(Y_1\), \(Y_2^\ast\), \(Y_4\) og \(Y_5\): likninger,
estimeringsprosedyre, kryssvalideringsresultater og framskrivningsmetode.
For \textbf{begrunnelsene} bak disse valgene (hvorfor Poisson, hvorfor
intercept-only, hvorfor 2010 er utelatt for \(Y_1\) osv.), se
\texttt{beslutningslogg.md}.

\textbf{NB}: Dette dokumentet er bevisst \textbf{ikke} tilgjengelig fra selve
Shiny-appen (\texttt{app.R}) -- se punkt 6 i beslutningsloggen.

\section{Datagrunnlag}

\begin{table}[htbp]
\centering
\small
\begin{tabularx}{\textwidth}{
    >{\raggedright\arraybackslash}p{1.2cm}
    >{\raggedright\arraybackslash}X
    >{\raggedright\arraybackslash}p{3.3cm}
    >{\raggedleft\arraybackslash}p{2.1cm}
    >{\raggedleft\arraybackslash}p{1.6cm}}
\toprule
\textbf{Variabel} & \textbf{Definisjon} & \textbf{Estimeringsperiode} &
\textbf{Kommune-år (N)} & \textbf{Kommuner} \\
\midrule
\(Y_1\) &
Brukere av hjemmetjenester, i alt &
2007--2025, ekskl. 2010 &
6\,337 & 357 \\

\(Y_2^\ast\) &
\(Y_2\) (heldøgnsbolig) + \(Y_{3a}\) (institusjon, langtid), i alt &
2007--2025 &
5\,167 & 355 \\

\(Y_4\) &
Gjennomsnittlig timer/uke, hjemmetjenester &
2007--2025 &
6\,329 & 357 \\

\(Y_5\) &
Sykepleiere, årsverk i alt &
2015--2025 (mangler helt 2007--2014) &
3\,898 & 357 \\
\bottomrule
\end{tabularx}
\end{table}

Alle modeller bruker
\[
\texttt{demensandel}
=
\frac{\texttt{demens\_estimert}}{\texttt{folk\_ialt}}
\]
(se \texttt{legg\_til\_demens.R}) og en tilfeldig kommuneeffekt over
\texttt{kommunenr\_2024} (2024-kommunestruktur, se
\texttt{hent\_paneldata\_2007\_2025.R}).

Demensestimatet beregnes som:
\[
\widehat{\text{demens}}_{k,t}
=
\sum_{a,s}
\left[
N_{a,s,k,t}\,
p_{a,s}
\right],
\]
der \(a\) angir alder, \(s\) kjønn, \(N_{a,s,k,t}\) er befolkningen i
alders-/kjønnsgruppe \(a,s\) i kommune \(k\) og år \(t\), og \(p_{a,s}\)
er den tilhørende prevalensraten.

Med variabelnavnene fra koden kan dette skrives som
\[
\texttt{demens\_estimert(kommune, år)}
=
\sum_{\text{alder, kjønn}}
\left[
\texttt{befolkning(alder, kjønn, kommune, år)}
\times
\texttt{prevalensrate(alder, kjønn)}
\right].
\]

Prevalensratene er faste, aldersgruppe- og kjønnsspesifikke andeler
(se \texttt{dementia\_dic} i \texttt{legg\_til\_demens.R}).

\section{Modell \(Y_1\) -- hjemmetjenestebrukere
(\texttt{modell\_y1.R})}

\textbf{Type}: Poisson generalisert lineær blandet modell
(\texttt{lme4::glmer}, log-lenke), med befolkning som eksponering
(offset).

\subsection*{Likning (hovedmodell)}

\begin{align}
Y_{1,a,\mathrm{ialt},k,t}
&\sim \operatorname{Poisson}(\mu_{kt}), \\
\log(\mu_{kt})
&=
\log(\texttt{folk\_ialt}_{kt})
+\beta_0
+\sum_t \gamma_t\,\text{år}_t
+\beta_1\,\texttt{demensandel}_{kt}
+u_k, \\
u_k &\sim \mathcal{N}(0,\sigma_u^2).
\end{align}

\subsection*{Estimerte parametre}

Hovedmodell, alle tilgjengelige år unntatt 2010:

\begin{table}[htbp]
\centering
\begin{tabularx}{0.9\textwidth}{>{\raggedright\arraybackslash}X >{\raggedright\arraybackslash}p{5.2cm}}
\toprule
\textbf{Parameter} & \textbf{Estimat} \\
\midrule
Intercept (\(\beta_0\)) & \(-3{,}832\) \\
\texttt{demensandel} (\(\beta_1\)) & \(34{,}02\) \\
Årseffekter (\(\gamma_t\)) &
spenner fra \(-0{,}047\) (2025) til \(0{,}059\) (2009), relativt til 2007 \\
Kommune-tilfeldig-intercept, SD (\(\sigma_u\)) & \(0{,}188\) \\
\bottomrule
\end{tabularx}
\end{table}

\subsection*{Kryssvalidering}

Leave-one-year-out, alle år:

\begin{table}[htbp]
\centering
\begin{tabularx}{0.9\textwidth}{
    >{\raggedright\arraybackslash}X
    >{\raggedleft\arraybackslash}p{1.7cm}
    >{\raggedleft\arraybackslash}p{1.7cm}
    >{\raggedleft\arraybackslash}p{2.2cm}}
\toprule
\textbf{Modellvariant} & \textbf{RMSE} & \textbf{MAE} & \textbf{Korrelasjon} \\
\midrule
Kun tilfeldig intercept (brukt til framskrivning) & 104,9 & 51,4 & 0,9964 \\
Intercept + korrelert helning på \texttt{demensandel} & 68,0 & 34,2 & 0,9985 \\
\bottomrule
\end{tabularx}
\end{table}

Helningsvarianten gir 91 av 357 kommuner en negativ effektiv helning på
\texttt{demensandel} -- derfor brukes intercept-only-modellen til
framskrivning, selv om helningsvarianten scorer bedre på CV
(se beslutningslogg pkt. 3).

\textbf{2010 er utelatt} fra estimeringsdataene for denne variabelen
(se beslutningslogg pkt. 1).

\section{Modell \(Y_2^\ast\) -- heldøgns omsorg
(\texttt{modell\_y2\_stjerne.R})}

\textbf{Type}: samme struktur som \(Y_1\) (Poisson GLMM, log-lenke,
befolkning som offset).

\subsection*{Likning}

\begin{align}
Y_{2^\ast,\mathrm{ialt},k,t}
&\sim \operatorname{Poisson}(\mu_{kt}), \\
\log(\mu_{kt})
&=
\log(\texttt{folk\_ialt}_{kt})
+\beta_0
+\sum_t \gamma_t\,\text{år}_t
+\beta_1\,\texttt{demensandel}_{kt}
+u_k .
\end{align}

der
\[
\texttt{Y\_2\_stjerne\_ialt}
=
\texttt{Y\_2\_ialt}
+
\texttt{Y\_3\_a\_ialt}.
\]

\subsection*{Estimerte parametre}

Hovedmodell, alle år 2007--2025:

\begin{table}[htbp]
\centering
\begin{tabularx}{0.9\textwidth}{>{\raggedright\arraybackslash}X >{\raggedright\arraybackslash}p{5.2cm}}
\toprule
\textbf{Parameter} & \textbf{Estimat} \\
\midrule
Intercept (\(\beta_0\)) & \(-4{,}956\) \\
\texttt{demensandel} (\(\beta_1\)) & \(30{,}80\) \\
Årseffekter (\(\gamma_t\)) &
monotont fallende fra 2010 (\(0{,}003\)) til 2025 (\(-0{,}171\)),
relativt til 2009 \\
Kommune-tilfeldig-intercept, SD (\(\sigma_u\)) & \(0{,}260\) \\
\bottomrule
\end{tabularx}
\end{table}

\subsection*{Kryssvalidering}

Leave-one-year-out, alle år:

\begin{table}[htbp]
\centering
\begin{tabularx}{0.9\textwidth}{
    >{\raggedright\arraybackslash}X
    >{\raggedleft\arraybackslash}p{1.7cm}
    >{\raggedleft\arraybackslash}p{1.7cm}
    >{\raggedleft\arraybackslash}p{2.2cm}}
\toprule
\textbf{Modellvariant} & \textbf{RMSE} & \textbf{MAE} & \textbf{Korrelasjon} \\
\midrule
Kun tilfeldig intercept (brukt til framskrivning) & 35,7 & 19,2 & 0,9968 \\
Intercept + korrelert helning på \texttt{demensandel} & 46,4 & 20,7 & 0,9950 \\
\bottomrule
\end{tabularx}
\end{table}

Her er intercept-only-modellen faktisk \textbf{best} på kryssvalidering i
tillegg til å unngå fortegnsproblemet (95/355 kommuner med negativ
effektiv helning i den korrelerte helningsvarianten, 78/355 i den
ukorrelerte).

\textbf{2010 er ikke utelatt} for denne variabelen -- sjekket spesifikt og
ikke funnet noe tilsvarende avvik (se beslutningslogg pkt. 1).

\section{Modell \(Y_4\) -- timer/uke (\texttt{modell\_y4.R})}

\textbf{Type}: lineær blandet modell (\texttt{lme4::lmer},
identitetslenke) -- \textbf{ikke} Poisson, siden \(Y_4\) allerede er et
gjennomsnitt (timer/uke per bruker), ikke et antall. Ingen offset.

\subsection*{Likning}

\begin{align}
Y_{4,\mathrm{ialt},k,t}
&=
\beta_0
+\sum_t \gamma_t\,\text{år}_t
+\beta_1\,\texttt{demensandel}_{kt}
+u_k+\varepsilon_{kt}, \\
u_k &\sim \mathcal{N}(0,\sigma_u^2), \\
\varepsilon_{kt} &\sim \mathcal{N}(0,\sigma_\varepsilon^2).
\end{align}

\subsection*{Estimerte parametre}

Hovedmodell, alle år 2007--2025:

\begin{table}[htbp]
\centering
\begin{tabularx}{0.9\textwidth}{>{\raggedright\arraybackslash}X >{\raggedright\arraybackslash}p{5.2cm}}
\toprule
\textbf{Parameter} & \textbf{Estimat} \\
\midrule
Intercept (\(\beta_0\)) & \(4{,}093\) \\
\texttt{demensandel} (\(\beta_1\)) & \(-0{,}899\) \\
Årseffekter (\(\gamma_t\)) &
stigende fra 2009 (\(0{,}460\)) til 2022 (\(1{,}022\)),
noe fallende mot 2025 (\(0{,}937\)) \\
Kommune-tilfeldig-intercept, SD (\(\sigma_u\)) & \(2{,}105\) \\
Residual-SD (\(\sigma_\varepsilon\)) & \(1{,}774\) \\
\bottomrule
\end{tabularx}
\end{table}

\subsection*{Kryssvalidering}

Leave-one-year-out, alle år:

\begin{table}[htbp]
\centering
\begin{tabularx}{0.9\textwidth}{
    >{\raggedright\arraybackslash}X
    >{\raggedleft\arraybackslash}p{1.7cm}
    >{\raggedleft\arraybackslash}p{1.7cm}
    >{\raggedleft\arraybackslash}p{2.2cm}}
\toprule
\textbf{Modellvariant} & \textbf{RMSE} & \textbf{MAE} & \textbf{Korrelasjon} \\
\midrule
Kun tilfeldig intercept (brukt til framskrivning) & 1,77 & 1,20 & 0,736 \\
Intercept + korrelert helning på \texttt{demensandel} & 1,53 & 0,98 & 0,813 \\
\bottomrule
\end{tabularx}
\end{table}

Vesentlig svakere prediktiv treffsikkerhet enn de øvrige variablene
(forventet -- se \texttt{modell\_y4.R}s egen kommentar om at \(Y_4\) er et
gjennomsnitt beregnet over få brukere i mange, særlig små, kommuner).
Helningsvarianten har et enda mer utbredt fortegnsproblem enn de andre
variablene: 195/357 kommuner (korrelert) / 192/357 (ukorrelert) får en
negativ effektiv helning på \texttt{demensandel}.

\textbf{Forkastet variant}: \texttt{Y\_1\_a\_ialt} (antall
hjemmetjenestemottakere) ble testet som en ekstra, rå forklaringsvariabel
på oppfordring. Ga ingen meningsfull forbedring i CV-korrelasjon
(\(\sim 0{,}735\) med og uten), og fikk \texttt{demensandel}s koeffisient
til å bli enda mer negativ -- et tegn på at \texttt{Y\_1\_a\_ialt} i stor
grad fanger opp kommunestørrelse (kollinearitet), ikke en selvstendig
effekt. Ikke tatt i bruk. Koden for dette forsøket ligger fortsatt i
\texttt{modell\_y4.R} (kommentert inn i formlene), men brukes ikke i
\texttt{framskriv\_y4.R} (som ikke er skrevet ennå).

\section{Modell \(Y_5\) -- sykepleierårsverk (\texttt{modell\_y5.R})}

\textbf{Type}: lineær blandet modell (\texttt{lme4::lmer}) på
\textbf{logaritmen} av \(Y_5\), med befolkning som offset (koeffisient
tvunget til 1) -- samme multiplikative populasjonsskalering som
Poisson-modellene for \(Y_1/Y_2^\ast\), men med normalfordelte feil på
logskala i stedet for Poisson-varians, siden årsverk ikke er et heltall.

\subsection*{Likning}

\begin{align}
\log(Y_{5,\mathrm{ialt},k,t})
&=
\log(\texttt{folk\_ialt}_{kt})
+\beta_0
+\sum_t \gamma_t\,\text{år}_t
+\beta_1\,\texttt{demensandel}_{kt}
+u_k+\varepsilon_{kt}, \\
u_k &\sim \mathcal{N}(0,\sigma_u^2), \\
\varepsilon_{kt} &\sim \mathcal{N}(0,\sigma_\varepsilon^2).
\end{align}

\subsection*{Estimerte parametre}

Hovedmodell, alle år 2015--2025:

\begin{table}[htbp]
\centering
\begin{tabularx}{0.9\textwidth}{>{\raggedright\arraybackslash}X >{\raggedright\arraybackslash}p{5.2cm}}
\toprule
\textbf{Parameter} & \textbf{Estimat} \\
\midrule
Intercept (\(\beta_0\)) & \(-5{,}256\) \\
\texttt{demensandel} (\(\beta_1\)) & \(13{,}83\) \\
Årseffekter (\(\gamma_t\)) &
stigende fra 2016 (\(0{,}016\)) til 2020 (\(0{,}070\)),
fallende mot 2025 (\(0{,}031\)), relativt til 2015 \\
Kommune-tilfeldig-intercept, SD (\(\sigma_u\)) & \(0{,}264\) \\
Residual-SD (\(\sigma_\varepsilon\)) & \(0{,}107\) \\
\bottomrule
\end{tabularx}
\end{table}

\subsection*{Kryssvalidering}

Leave-one-year-out, alle 11 år:

\begin{table}[htbp]
\centering
\begin{tabularx}{0.9\textwidth}{
    >{\raggedright\arraybackslash}X
    >{\raggedleft\arraybackslash}p{1.7cm}
    >{\raggedleft\arraybackslash}p{1.7cm}
    >{\raggedleft\arraybackslash}p{2.2cm}}
\toprule
\textbf{Modellvariant} & \textbf{RMSE} & \textbf{MAE} & \textbf{Korrelasjon} \\
\midrule
Kun tilfeldig intercept (brukt til framskrivning) & 13,9 & 5,6 & 0,9975 \\
Intercept + korrelert helning på \texttt{demensandel} & 12,7 & 4,9 & 0,9979 \\
\bottomrule
\end{tabularx}
\end{table}

God prediktiv treffsikkerhet, på linje med \(Y_1/Y_2^\ast\).
Helningsvarianten: 103/357 kommuner med negativ effektiv helning
(korrelert), men bare 10/357 i den \textbf{ukorrelerte} varianten --
vesentlig bedre enn for de andre variablene, men intercept-only brukes
likevel til framskrivning, i tråd med prosjektets generelle policy.

\textbf{Datadekning}: kun 2015--2025 (11 år) har faktiske observasjoner --
2007--2014 mangler helt i datagrunnlaget. Modellen er derfor trent på et
vesentlig kortere tidsvindu enn de andre variablene.

\section{Kryssvalideringsmetode (gjelder alle modeller)}

Leave-one-year-out (LOYO): for hvert år \(t\) i datasettet estimeres
modellen på \textbf{alle andre år}, og brukes til å predikere år \(t\).
Siden alle modellene har årsdummyer (\texttt{år\_f}), har den
re-estimerte modellen per definisjon ingen egen koeffisient for det
utelatte året. Dette løses med en forenkling:
\[
\gamma_t^{(\text{utelatt år})}
\approx
\operatorname{gjennomsnitt}
\left(
\gamma_s:\; s \neq t
\right).
\]

Dette er en tilnærming, ikke en eksakt løsning -- se
\texttt{loyo\_cv()}-funksjonen i hvert \texttt{modell\_*.R}-script.
RMSE, MAE og korrelasjon rapporteres både per utelatt år og samlet over
alle utelatte år.

\section{Framskrivningsmetode
(\texttt{framskriv\_y1.R}, \texttt{framskriv\_y2\_stjerne.R},
\texttt{framskriv\_y5.R})}

\textbf{Oppdatert metode (rettet policy -- se beslutningslogg.md,
``Retur til slope-modell'')}: Punktestimatet bruker nå
\texttt{hovedmodell\_slope} (tilfeldig intercept \textbf{og} helning på
\texttt{demensandel} per kommune), \textbf{ankret} ved kommunens siste
observerte (2025) demensandel -- \textbf{ikke} en separat
intercept-only-modell som tidligere. Se punkt 2 i denne oppdateringen for
hvorfor.

\subsection*{Trinn}

\begin{enumerate}[leftmargin=*,label=\arabic*.]
\item Hent SSBs befolkningsframskrivning (tabell 12882, hovedalternativ
MMMM = \texttt{ContentsCode = "Personer"}) for 2026--2050, per kommune,
alder og kjønn.

\item Beregn framskrevet \texttt{demensandel} med \textbf{samme}
prevalensrater som i historikken.

\item Predikér med \texttt{hovedmodell\_slope}, \textbf{ankret} ved
kommunens siste observerte (2025) demensandel (\(a_k\)). Årseffekten for
et framtidig år settes til gjennomsnittet av de historiske årseffektene,
som før. La \(u_{k0}\) og \(u_{k1}\) være kommunens tilfeldige intercept
og helning:
\begin{align}
\mathrm{intercept}'_k
&=
\beta_0+u_{k0}+(\beta_1+u_{k1})a_k,
\\
\widehat{y}_{kt}
&=
\exp\left[
\mathrm{intercept}'_k
+\overline{\gamma}
+\beta_1\bigl(
\texttt{demensandel}_{kt}-a_k
\bigr)
+\log(\texttt{folk\_ialt}_{kt})
\right].
\end{align}

\textbf{Hvorfor ankring, og hvorfor dette gjør slope-modellen trygg å
bruke}: \(\mathrm{intercept}'_k\) er en \textbf{konstant} (bygget fra
modellens fit ved ankeret) og kan derfor ikke selv eksplodere. All
\textbf{framtidig} vekst i demensandel bruker bare det faste, nasjonale
helningsanslaget \(\beta_1\) -- aldri kommunens egen tilfeldige helning
\(u_{k1}\). Det er nettopp bruken av \(u_{k1}\) på framtidig vekst som
tidligere ga et implausibelt fortegn for en andel kommuner (se pkt. 3 i
hoveddokumentet). Siden \(u_{k1}\) her bare multipliserer et fast
historisk ankerpunkt (ikke en økende framtidig størrelse), er denne faren
eliminert.

Uten ankring -- dvs. å bruke
\(u_{k1}\cdot\texttt{demensandel}_{kt}\) direkte, eller å bytte mellom en
egen intercept-only-modell og denne modellen -- oppstår et reelt
inkonsistensproblem: de to modellvariantene har ulike faste effekter
(\(\beta_0,\beta_1\)). Et brukergrensesnitt som lar brukeren ``slå av''
en tilfeldig-helning-effekt ved å bytte til en annen modell vil derfor
vise et synlig sprang selv når effekten skulle vært 0\,\%. Dette ble
funnet som en reell bug i appen og rettet ved å alltid bruke
\texttt{hovedmodell\_slope}, aldri en egen intercept-only-modell, for
framskrivning.

\item \textbf{Glidende overgang} fra observert 2025-nivå til
modellframskrivning, for å unngå brå hopp i grafene ved skjøten mellom
historikk og framskrivning. Overgangen er nå strukket over
\textbf{hele framskrivningsperioden} (2026--2050) i stedet for bare de
første 10 årene, for å unngå et kink i grafen der overgangsvekten
tidligere flatet ut brått til 0 i 2036:
\begin{align}
w_t
&=
\max\left(
0,\;
0{,}9\frac{2050-t}{2050-2026}
\right),
\qquad t=2026,\ldots,2050,
\\
y_{\mathrm{predikert},t}
&=
w_t\,y_{\mathrm{observert},2025}
+
(1-w_t)\widehat{y}_{kt}.
\end{align}

Dette gir 90\,\% vekt på observert 2025-nivå i 2026, lineært avtagende
til 0\,\% (100\,\% modell) først i 2050 (siste framskrevne år). Se
\texttt{glatt\_overgang()} i hvert \texttt{framskriv\_*.R}-script.
NB: Dette betyr at observert 2025-nivå har en god del gjenværende vekt
langt inn i perioden (f.eks. \(\sim 45\)--50\,\% ved 2036) -- en bevisst
avveining mellom å unngå kink-artefaktet og hvor lenge ett enkelt
historisk datapunkt påvirker framskrivningen.
\end{enumerate}

\subsection*{Kommunespesifikk trend (appen)}

\texttt{app.R} sin \texttt{framskriv\_trend()} bruker nøyaktig samme
ankermetode som over, men med en brukerstyrt vekt
\[
\texttt{vekt\_trend}(t)
=
\max\left(
0,\;
\frac{T-(t-2026)}{T}
\right)
\]
multiplisert med \(u_{k1}\), i stedet for at helningsbidraget alltid er 0.
Dette er en generalisering av pkt. 3; ``standardvisningen'' i appen er
nøyaktig det samme som \(T=0\). Se kommentarblokken ved
\texttt{framskriv\_trend()} i \texttt{app.R}.

\subsection*{Befolkningsscenario}

Kun ett befolkningsscenario (MMMM) er beregnet så langt.
\texttt{ALTERNATIV}-variabelen i skriptene kan endres til
\texttt{"Personer1"} (LLML) eller \texttt{"Personer2"} (HHMH) for å
beregne alternative scenarioer.

\section{Kjente begrensninger / videre arbeid}

\begin{itemize}[leftmargin=*]
\item Ingen konfidens- eller prediksjonsintervall er beregnet ennå
(bare punktestimater).
\item Kun ett befolkningsscenario (MMMM) er beregnet.
\item \(Y_3\) er ikke modellert som egen variabel (kun som del av
\(Y_2^\ast\)).
\item \texttt{framskriv\_y4.R} er ikke skrevet --
\(Y_4\)s svake prediktive treffsikkerhet (korrelasjon \(\sim 0{,}74\))
gjør at videre arbeid her bør vurderes nøye før framskrivning tas i bruk.
\item LOYO-kryssvalideringens håndtering av årseffekt for utelatt år
(gjennomsnitt av andre år) er en forenkling -- se pkt. 6.
\end{itemize}

\end{document}
